\begin{tikzpicture}
\node[note,text width=13cm] at (4,4.6) {共同初态：j=40，余额 10000，M/D 为空；下一条为 41 的 S1};
\node[warn,text width=5.4cm] (a) at (0,2.9) {先保存业务，再保存边界\\余额 7000、D1 已在\\M 漏写，j 仍为 40 时崩溃};
\node[warn,text width=5.4cm] (b) at (8,2.9) {先保存边界，再保存业务\\j 已为 41\\余额仍为 10000 时崩溃};
\node[warn,text width=5.4cm] (aa) at (0,.8) {重放 41\\M 为空，可能再次减 3000\\余额成为 4000};
\node[warn,text width=5.4cm] (bb) at (8,.8) {从 42 继续\\跳过尚未体现在业务中的 S1\\余额错误保持 10000};
\draw[arr] (a)--(aa);\draw[arr] (b)--(bb);
\node[good,text width=13cm] at (4,-1.3) {把 W、M 与 j 放入同一本地事务\\崩溃后只会见到完整的 40 状态，或完整的 41 状态};
\end{tikzpicture}
